\documentclass[11pt]{article}

\usepackage[utf8]{inputenc}
\usepackage[T1]{fontenc}
\usepackage[a4paper,margin=1in]{geometry}
\usepackage{amsmath,amssymb}
\usepackage{graphicx}
\usepackage{booktabs}
\usepackage{siunitx}
\usepackage{xcolor}
\usepackage[colorlinks,linkcolor=blue,citecolor=blue,urlcolor=blue]{hyperref}

\graphicspath{{figures/}}
\sisetup{detect-all, range-phrase={--}, range-units=single}
\setlength{\parskip}{0.4em}
\setlength{\parindent}{0pt}

\title{\vspace{-2em}Distance and uncertainty for the normalized SAR values}
\author{Robin Wydaeghe}
\date{For Hamed and the dose-model team \quad\textbar\quad June 2026}

\begin{document}
\maketitle
\vspace{-1.5em}

You asked for two things at the dose-model meeting. First, a way to correct a
normalized SAR value when the phone is at a different distance. Second, an
uncertainty estimate for each value, such as the \SI{2.62}{W\per kg\per W} brain
figure for Duke in the front-of-eyes scenario. Both are ready. The numbers and
tables are in the attached Excel file \texttt{nearfield\_dose\_results.xlsx}, and
the figures below are also in the email.

Each scenario now has a full distance curve and an uncertainty range, not a single
number.

\section*{The method}

This is my method. It is a fast surface model: it computes the dose on the
\SI{}{2D} body surface, not in the full \SI{}{3D} tissue volume. A full-wave
(FDTD) run gives one number for one phone position and orientation, and takes
hours. Mine gives thousands of positions and orientations in minutes. It does not
replace full-wave results.

Figure~\ref{fig:geometry} shows the three phone placements on the Duke model.
Figure~\ref{fig:apd} shows the absorbed power on the skin. The dose is highest on
the face nearest the phone.

\begin{figure}[h]
  \centering
  \includegraphics[width=\linewidth]{fig_geometry.pdf}
  \caption{The three hand-held placements on Duke: in front of the eyes
  (\SI{200}{mm}), at the cheek (\SI{8}{mm}), and in front of the belly
  (\SI{200}{mm}). The face is at $-y$.}
  \label{fig:geometry}
\end{figure}

\begin{figure}[h]
  \centering
  \includegraphics[width=\linewidth]{fig_apd_map.pdf}
  \caption{Absorbed power on the skin for the front-of-eyes case at
  \SI{2450}{MHz}, per watt radiated. Front view and side view. The dose is highest
  on the near cheek and brow.}
  \label{fig:apd}
\end{figure}

\section*{1. Distance correction}

The dose near the body is inverse-square in distance, with a small offset:
\begin{equation}
  \label{eq:dist}
  \mathrm{SAR}(d) = \mathrm{SAR}(d_{\mathrm{ref}})\,
  \left(\frac{d_{\mathrm{ref}}+\delta}{d+\delta}\right)^{2}.
\end{equation}
Here $d$ is the phone-to-body distance, $d_{\mathrm{ref}}$ is the distance the
value was computed at (\SI{200}{mm} for front of eyes), and $\delta$ is a small
near-field offset of about \SI{6}{mm}. This is the textbook inverse-square rule
beyond a few centimetres. Halve the distance and the dose is about four times
higher. The fit is good: $R^2 > 0.999$ for every band and placement.

The brain value of \SI{2.62}{W\per kg\per W} at \SI{200}{mm} becomes:

\begin{table}[h]
  \centering
  \caption{The \SI{2.62}{W\per kg\per W} brain value at other distances, from
  Eq.~\eqref{eq:dist}. The last column is the $\pm1\sigma$ orientation band
  (Section~2).}
  \label{tab:applied}
  \begin{tabular}{rccc}
    \toprule
    $d$ (\si{mm}) & factor & adjusted (\si{W\per kg\per W}) & $\pm1\sigma$ \\
    \midrule
    100 & 3.76 & 9.86 & 7.2--12.5 \\
    150 & 1.74 & 4.56 & 3.3--5.8 \\
    200 & 1.00 & 2.62 & 1.9--3.3 \\
    250 & 0.65 & 1.70 & 1.2--2.2 \\
    300 & 0.45 & 1.19 & 0.9--1.5 \\
    400 & 0.26 & 0.68 & 0.5--0.9 \\
    \bottomrule
  \end{tabular}
\end{table}

Figure~\ref{fig:factor} plots this, with the reference point marked.
Figure~\ref{fig:bands} shows that the law holds for all eight bands. The per-band
fit coefficients are in the Excel file (sheets \texttt{distlaw\_*}).

\begin{figure}[h]
  \centering
  \includegraphics[width=\linewidth]{fig_distance_factor.pdf}
  \caption{The brain value adjusted for distance, with its orientation band. The
  red dot is the \SI{2.62}{W\per kg\per W} reference at \SI{200}{mm}.}
  \label{fig:factor}
\end{figure}

\begin{figure}[h]
  \centering
  \includegraphics[width=\linewidth]{fig_distance_bands.pdf}
  \caption{The distance law for all eight bands (front of eyes, Duke). The dotted
  \SI{700}{} and \SI{835}{MHz} curves are flagged: their shape is correct, their
  absolute level is not (see caveats).}
  \label{fig:bands}
\end{figure}

\section*{2. Uncertainty}

I give the mean, standard deviation, min, max, and the \SIrange{5}{95}{}
percentiles for each value. For a fixed body the main source is the phone
orientation. I sample it with a large set of random orientations.

The orientation spread for the Duke brain value is about \SI{27}{\percent} (one
standard deviation). So
\[
  \SI{2.62}{W\per kg\per W} \text{ becomes } 2.6 \pm 0.7~\si{W\per kg\per W},
\]
and over all orientations the full range is \SIrange{0.8}{3.8}{W\per kg\per W}. The
straight-at-the-head pose is not the worst case, so the spread is not symmetric.

The body model is a second source. The value differs by a further
\SIrange{30}{40}{\percent} across the four models (Duke, Ella, Eartha,
Thelonious). Smaller
bodies absorb more per kilogram, so a child model is higher than the adult. The two
sources add in quadrature, for about \SI{45}{\percent} in total.
Figure~\ref{fig:unc} shows the orientation spread per placement. The per-scenario
tables are in the Excel file (sheets \texttt{unc\_*}).

My method is a third source, and the largest one on absolute values. Against
full-wave it is within a few percent near the top of the range, and it
underestimates toward the lowest bands, by up to roughly a factor of two at
\SI{700}{MHz}. This does not reach the numbers here. The absolute level is your
full-wave value, and my method gives only the distance correction and the spread.
Both are ratios, and a steady per-band bias cancels in a ratio. It would matter
only for an absolute SAR read straight from the method, which the dose model does
not use.

\begin{figure}[h]
  \centering
  \includegraphics[width=\linewidth]{fig_uncertainty.pdf}
  \caption{Orientation spread at the nominal distance, per placement
  (\SI{2450}{MHz}, Duke), each normalized to its straight-at-the-body pose.}
  \label{fig:unc}
\end{figure}

\section*{The Excel file}

The Excel file has:
\begin{itemize}\setlength{\itemsep}{0pt}
  \item \texttt{README}: the sheet descriptions and the normalization.
  \item \texttt{distlaw\_*}: the distance-law coefficients ($\delta$, $R^2$) for
  every band and metric.
  \item \texttt{unc\_*}: mean, std, min, max, and percentiles per scenario.
  \item \texttt{duke\_fronteyes\_detail}: per-distance statistics for the
  front-of-eyes case.
  \item \texttt{brain\_2.62\_adjusted}: Table~\ref{tab:applied} in full.
\end{itemize}

\section*{Caveats}

The brain number is a volume quantity from your full-wave model. I do not
recompute it. I take that value and report how it scales with distance and how
much it varies.

The near-field placement adds one effect the full-wave validation does not cover:
the coupling between the phone and the body at the smallest gaps. So the absolute
SAR is least certain at the cheek and the lowest bands. Beyond a few centimetres
the scaling rests on the geometry, and the closest distances (under a couple of
centimetres) are an upper bound.

The \SI{700}{} and \SI{835}{MHz} antennas are off-resonance in free space, so their
absolute level is unreliable. Their shape with distance and orientation is still
correct, and the scaling and the spread depend only on that shape.

I can extend this to other tissues, bands, or scenarios.

\end{document}
